\documentclass[11pt]{article}
%==============================================================================
%  Reproducing Our Table 3: A Campaign Runbook
%  How the earth_elliptic_to_geo min-fuel thrust ladder was actually solved,
%  per rung, with the reasoning -- and why run_gergaud will not reproduce the
%  deep ladder.
%
%  Source of record for numbers/provenance:
%    earth_elliptic_to_geo/CAMPAIGN.md
%    earth_elliptic_to_geo/.superpowers/sdd/progress.md  (MEE ladder ledger)
%    earth_elliptic_to_geo/doc/table3_method_note.tex   (the method/theory)
%==============================================================================
\usepackage[margin=1in]{geometry}
\usepackage{amsmath,amssymb}
\usepackage{booktabs}
\usepackage{xcolor}
\usepackage{listings}
\usepackage[colorlinks=true,linkcolor=blue,citecolor=blue,urlcolor=blue]{hyperref}
\usepackage{enumitem}
\setlist{itemsep=2pt,topsep=3pt}
\emergencystretch=2.5em

\definecolor{codegray}{gray}{0.95}
\definecolor{codekw}{rgb}{0.13,0.13,0.7}
\definecolor{codecomment}{rgb}{0.0,0.5,0.0}
\definecolor{codestr}{rgb}{0.6,0.1,0.1}
\lstdefinestyle{mat}{
  language=Matlab, basicstyle=\ttfamily\footnotesize,
  keywordstyle=\color{codekw}, commentstyle=\color{codecomment}\itshape,
  stringstyle=\color{codestr}, backgroundcolor=\color{codegray},
  numbers=none, breaklines=true, frame=single, framerule=0pt,
  showstringspaces=false, columns=fullflexible, keepspaces=true, aboveskip=6pt, belowskip=6pt}
\lstset{style=mat}
\newcommand{\code}[1]{\texttt{\small #1}}
\newcommand{\file}[1]{\texttt{\small #1}}
\newcommand{\dL}{\Delta L}
\newcommand{\Tmax}{T_{\max}}

\title{\vspace{-1.5em}Reproducing Our Table~3: A Campaign Runbook\\[3pt]
\large How the minimum-fuel thrust ladder was actually solved, per rung ---
and why \file{run\_gergaud} will not reproduce the deep ladder}
\author{\normalsize \file{earth\_elliptic\_to\_geo} pipeline --- internal runbook}
\date{\normalsize 2026-07-18}

\begin{document}
\maketitle
\vspace{-1.4em}

\begin{abstract}
\noindent
This runbook records \emph{exactly} how each row of our version of Gergaud's
Table~3 was produced during the campaign, and the reasoning behind every
non-obvious step. It is a companion to the theory note
(\file{doc/table3\_method\_note.tex}); read that first for the optimal-control
problem, the MEE formulation, and the discretization. The short version: the
ladder was a \emph{warm-chained thrust continuation} with \emph{per-rung
recipes} that get progressively more hands-on as thrust drops, and the newer
front door \file{run\_gergaud.m} --- built after the campaign --- deliberately
does \emph{not} carry those recipes, so it is the wrong tool for re-solving the
hard rungs (\S\ref{sec:whynot}).
\end{abstract}

\section{Two meanings of ``reproduce''}

\paragraph{(A) Present the numbers we already certified.} If the goal is the
table/figure as reported, nothing needs re-solving --- the certified solutions
are cached in \file{results/} (Table~\ref{tab:provenance}). Use the campaign's
own generators, which read those caches:
\begin{lstlisting}
run_ladder([10 5 2.5 1])   % prints the ladder summary table from the rung caches
fig_table3                 % the Table-3 figure (switches / revs / R0 panels)
\end{lstlisting}
\file{run\_gergaud(struct('thrustN',T,'runMode','auto'))} returns the same
numbers one row at a time (and can add a plot/movie); it is the right tool for a
single row or a custom-endpoint row, not for regenerating the whole table.

\paragraph{(B) Re-solve from scratch.} This is the rest of the document: the
per-rung recipes (\S\ref{sec:rungs}), the operational essentials that keep them
converging (\S\ref{sec:ops}), and why \file{run\_gergaud} cannot stand in for
them (\S\ref{sec:whynot}).

\begin{table}[h]\centering\small
\begin{tabular}{@{}llll@{}}
\toprule
$\Tmax$ & fuel result (\file{results/}) & min-time anchor & headline $m_f$ / sw / revs \\
\midrule
10\,N   & \file{MEE\_M2\_10N.mat}                 & \file{MEE\_mintime\_T100.mat} & 1377.10 / 19 / 7.33 \\
5\,N    & \file{MEE\_M2\_5N.mat}                  & \file{MEE\_mintime\_T50.mat}  & 1364.54 / 32 / 14.16 \\
2.5\,N  & \file{MEE\_M2\_2p5N.mat}                & \file{MEE\_mintime\_T25.mat}  & 1369.79 / 76 / 27.84 \\
1\,N    & \file{MEE\_M2\_1N\_PSR\_psr\_final.mat} & \file{MEE\_mintime\_T10.mat}  & 1371.44 / 171 / 69.15 \\
0.5\,N  & \file{MEE\_M2\_0p5N\_PSR\_psr\_final.mat} & \emph{none (R0-law)}       & 1375.28 / 362 / 138.60 \\
\bottomrule
\end{tabular}
\caption{Where each certified row lives. The 1\,N and 0.5\,N headline rows are
the PSR-refined files (they supersede the coarse \file{MEE\_M2\_1N}/\file{\_0p5N}).
The 0.5\,N anchor is an R0-law estimate, not a certified min-time solve.}
\label{tab:provenance}
\end{table}

\section{The strategy in one paragraph}

Each row is a fixed-time min-fuel solve at $t_f=c_{tf}\,t_{f,\min}(\Tmax)$ with
$c_{tf}=1.5$, so every rung needs (i) a min-time anchor $t_{f,\min}$ and (ii) an
energy$\to$fuel homotopy at that fixed time. The ladder descends
$10\to5\to2.5\to1\to0.5$~N, \emph{warm-chaining} each rung from the one above:
because $\Tmax\,t_{f,\min}\approx\text{const}$ (the $R0$ law), the total winding
$\dL\sim1/\Tmax$ at fixed trajectory shape, so the previous rung's converged
trajectory --- rescaled $\dL_{\text{new}}=\dL_{\text{prev}}(\Tmax^{\text{prev}}/
\Tmax^{\text{new}})$ --- is an excellent starting shape for both the anchor and
the fuel solve. This warm chain is what \file{run\_ladder.m} automates. The
recipes below are what that chain needs at each rung to actually converge.

\section{Per-rung runbook}
\label{sec:rungs}

\subsection*{10\,N --- the anchor of the ladder}
\begin{lstlisting}
% fuel solve (energy->fuel homotopy at fixed t_f):
res = run_transfer_mee(struct('thrustN',10,'ctf',1.5,'tfMinAnchor',22.2206));
% min-time anchor (multi-basin -- keep-best of cold + warm):
mt  = run_mintime_mee(10, 25);
\end{lstlisting}
\textbf{Why / what to watch:}
\begin{itemize}
\item \textbf{Seed throttle $\approx0.4$, not 1.} A full-throttle constant burn
from the $e_0=0.75$ apogee overshoots GEO by revolution $\sim3$ and goes
coordinate-singular; \file{mee\_seed} uses $\delta\!\approx\!0.4$, \code{stopP=1},
giving a $\sim7.3$-rev, well-conditioned seed.
\item \textbf{Node density $25$/rev is production.} A warm mesh-refinement study
(15/20/25/30/40 nodes/rev) showed $m_f$ flat to $0.025$~kg and \code{sw=19}
across $\{25,30,40\}$; \textbf{cold} solves at $\ge30$/rev are ill-conditioned
(dual infeasibility $\sim10^{16}$), so always \emph{warm}-refine dense meshes,
never cold-solve them.
\item \textbf{Min-time is multi-basin.} \file{run\_mintime\_mee} runs both a
fuel-warm Stage~A \emph{and} a cold 3-rev tangential Stage~B and keeps the lower
$t_f$. Here Stage~A converged \emph{cleanly} to a spurious $7.16$-rev basin
($14\%$ off); the cold Stage~B found the correct $4.503$-rev basin
($t_{f,\min}=22.2206$, matching the Cartesian anchor's $4.504$ revs). ``Converged''
is not evidence of the globally relevant optimum.
\end{itemize}

\subsection*{5\,N and 2.5\,N --- the warm-chained middle}
\begin{lstlisting}
run_ladder([10 5 2.5], struct('mtMaxIter', 300));
\end{lstlisting}
\textbf{Why / what to watch:}
\begin{itemize}
\item \textbf{5\,N broke the Cartesian wall.} The old Cartesian/Sundman stack
died at 5\,N; the MEE$+\dL$ formulation, warm-chained from 10\,N, certifies it
cleanly ($t_{f,\min}=44.6796$, $m_f=1364.54$, \code{sw=32}).
\item \textbf{2.5\,N ``wall'' was a budget artifact, not conditioning.} The
first 2.5\,N attempts stalled; a probe proved fresh solves converge to $\sim
10^{-14}$ in $<110$ iterations. The cause was a per-round iteration cap
(\code{mtMaxIter=50}) combined with IPOPT's barrier parameter resetting each
round --- every capped round re-fought the same barrier transition. \textbf{Cure:
raise the per-round budget to $\ge150$--$300$} (\code{cfg.mtMaxIter}). This is why
the call above passes \code{mtMaxIter=300}.
\end{itemize}

\subsection*{1\,N --- the hard anchor (small-N-first), plus a real infeasibility bug}
This rung does \emph{not} converge by simply asking for 25/rev. The recipe:
\begin{lstlisting}
% ANCHOR, step 1: small-N-first at 15 nodes/rev (converges where 25/rev crashes)
mt15 = run_mintime_mee(1, 15);          % ~6 min; manual continuation past the stall guard
% ANCHOR, step 2: warm mesh-refine 15/rev -> 25/rev in ONE tight solve (~57 s)
%   (interp_warmstart onto the 25/rev grid, then a single warmTight casadi_lt_mee)
% FUEL: at the certified anchor, energy->fuel homotopy
res1 = run_transfer_mee(struct('thrustN',1,'ctf',1.5,'tfMinAnchor',mt.tfmin));
% PSR: switch-time sharpening (headline result is round 2)
psr  = psr_mee_refine(res1, struct( ... ));   % -> MEE_M2_1N_PSR_psr_final.mat, 1371.44 kg
\end{lstlisting}
\textbf{Why / what to watch:}
\begin{itemize}
\item \textbf{Small-N-first, then refine.} A direct 25/rev ($N\!\approx\!1104$)
cold anchor hit four reproducible mid-computation MEX/SIGBUS crashes; the same
$N=1104$ solved in $57$~s \emph{as a warm refinement} from a converged 15/rev
($N=662$) anchor. The crashes were conditioning-driven, not size-driven ---
approach the hard size warm, never cold.
\item \textbf{The fuel solve had a structural-infeasibility bug.} A hardcoded
time-box $t\le300$ collided with $t_f=335.7$: the fixed-$t_f$ boundary condition
and the box were mutually unsatisfiable before IPOPT iterated (IPOPT's ``NLP
error $35.46$'' $=335.71-300.25$ was the tell; \code{maxDefect}/\code{termErr}
never checked the time equality). \textbf{Fix: the time box is now
tfTarget-relative} --- already in \file{casadi\_lt\_mee.m}. With that, 1\,N fuel
certified on the first attempt ($m_f=1370.36$, \code{sw=171}).
\item \textbf{PSR gives the headline.} \file{psr\_mee\_refine} round~2 sharpens
switch times, moving $m_f$ $1370.36\to1371.44$ (edge $0.9983\to0.9994$, revs
stable) --- a discretization refinement, not a basin change. The $171$-vs-$179$
switch-count gap versus the paper is open (windowed PSR scoring cannot
\emph{discover} new switch pairs; a periodic uniform-sweep round is the fix).
\end{itemize}

\subsection*{0.5\,N --- anchor-free (the min-time anchor is walled)}
\begin{lstlisting}
% NO certified min-time anchor exists. Use the R0-law estimate:
%   t_f = c_tf * (R0 / T),  R0 = mean(T*t_fmin) = 223.14 ND  (from the 4 certified anchors)
%   t_f = 1.5 * (223.14 / 0.5) = 669.42 ND
res05 = run_transfer_mee(struct('thrustN',0.5,'ctf',1.5,'tfMinAnchor',223.14/0.5));
psr05 = psr_mee_refine(res05, struct( ... ));   % 5 rounds (hybrid global) -> 1375.28 kg, sw=362
\end{lstlisting}
\textbf{Why / what to watch:}
\begin{itemize}
\item \textbf{The min-time anchor hit a conditioning wall.} Seven configurations
were tried; the best defect was $0.0545$ ($12$/rev $\times75$; $15$/rev was
worse), with a bit-identical stall on retry and repeated \code{libcoinmumps}
MEX/SIGBUS crashes. So the 0.5\,N row is built against an \emph{R0-law estimate}
$t_{f,\min}\approx446.27$~ND, not a certified anchor.
\item \textbf{The R0 law makes this honest but caveated.} $\Tmax\,t_{f,\min}
\approx850$~N$\cdot$h ($223.14$~ND) holds to $0.72\%$ across the four certified
anchors, so the estimate is well grounded --- but the 0.5\,N $t_f/t_{f,\min}$
column is a construction, \emph{excluded} from any $R0$ refit (circularity), and
the $m_f=1375.28$/\code{sw}$=362$ result is PSR-round-5-of-5, budget-limited
(\code{stopReason=maxRounds}), not confirmed mesh-stable.
\end{itemize}

\subsection*{0.2\,N and 0.1\,N --- not attained}
Honestly never solved: the deep-descent effort stopped at the 0.5\,N min-time
wall. Reaching them means either certifying a 0.5\,N anchor and chaining down, or
extending the anchor-free R0-law path to $\sim300$ and $\sim600$ switches at $N$
in the tens of thousands --- expecting the same crash class. See \file{TODO.md}
items~1--2.

\section{Operational essentials (the whole ladder)}
\label{sec:ops}
\begin{enumerate}[leftmargin=1.4em]
\item \textbf{Warm-start dense meshes; never cold-solve them.} $\ge30$/rev cold
is ill-conditioned. Mesh-refine from a certified coarser solve
(\file{interp\_warmstart}: linear for the state and RTN direction, nearest for
the throttle to keep switch edges crisp, then one \code{warmTight} solve).
\item \textbf{Per-round iteration budget $\ge150$--$300$.} A small cap plus
IPOPT's per-round barrier reset stalls the middle rungs (the false 2.5\,N wall).
\item \textbf{tfTarget-relative time box.} Never a hardcoded $t\le300$; derive it
from $t_f$ (the 1\,N infeasibility bug).
\item \textbf{\code{mumps\_pivot\_order=0}.} METIS ordering hard-aborts at
$N\!\sim\!193$ on this machine; AMD ordering fixes it.
\item \textbf{Never cache an uncertified result.} A loose iterate must not poison
the warm-start chain. Every driver caches only defect/terminal-certified
solutions; per-round and per-$\varepsilon$-step caches make a killed run
resumable by simply re-invoking.
\item \textbf{Classify crashes by the dump, not the exit code.} The recurring
\code{libcoinmumps} MEX/SIGBUS death is volume/duration-linked; a fresh
\file{\textasciitilde/matlab\_crash\_dump.*} identifies it (exit $137$ is
non-discriminating). Recovery is a plain relaunch (the caches resume).
\item \textbf{The $R0$ law is the anchor-free fallback.} $\Tmax\,t_{f,\min}
\approx223.14$~ND supplies a target when a min-time anchor will not certify
(the 0.5\,N recipe).
\end{enumerate}

\section{Why \file{run\_gergaud} will not reproduce the deep ladder}
\label{sec:whynot}

\file{run\_gergaud} is the right front door for a \emph{single} row, a
\emph{custom-endpoint} row, or a plot/movie. It is the wrong tool for re-solving
(B) the low-thrust rungs, for four concrete reasons:
\begin{enumerate}[leftmargin=1.4em]
\item \textbf{It is single-rung; the campaign was warm-chained.} The
$5\to2.5\to1$~N descent worked because each rung warm-started from the one above
(the $\dL$ $R0$-rescale in \file{run\_ladder}). \file{run\_gergaud solve} treats
one thrust in isolation, so a 1\,N solve starts far colder and is much harder.
\item \textbf{It does not encode the per-rung recipes.} \file{run\_gergaud solve}
calls \file{run\_mintime\_mee($T$,25)} directly --- straight into the 25/rev
crash wall at 1\,N (it has no ``small-N-first then refine'' step) --- and has no
notion of the 0.5\,N anchor-free R0-law construction or its multi-round PSR.
\item \textbf{Its solve/probe path was wired but never run to convergence on the
deep rungs.} The build's final review checked it as \emph{correct wiring}, not as
a validated deep solve; the certified numbers came from the campaign's driver
runs, not from \file{run\_gergaud}.
\item \textbf{For 0.5\,N its \code{auto} row already carries an R0-law $t_{f,\min}$
estimate (446.27\,ND) as a footnote} --- correct for \emph{presenting} the row,
but re-deriving it would not reconstruct the 5-round, budget-limited PSR process
that actually produced $m_f=1375.28$.
\end{enumerate}
In short: use \file{run\_ladder} + the recipes above (and the ledger/reports for
the exact hand-steps) to re-solve; use \file{run\_gergaud} to present a row, to
explore a custom transfer, or to render a movie.

\end{document}
